\section{Normalize}

\paragraph{Input and output.}
Normalize consumes a validated closed AST and emits a normalized closed AST.

\paragraph{Contract.}
The stage removes call and control sugar that would otherwise complicate
flattening.  Callee path suffixes become ordinary \code{Path}-typed arguments,
pipe syntax becomes ordinary calls, callable-name expressions that denote
zero-argument calls become explicit calls, and loop bodies receive explicit
yield/carry protocol structure.

\paragraph{Exclusions.}
Normalize does not globally expand default arguments.  Defaults are handled by
the elaboration stage and by type-aware binding where appropriate.

\paragraph{Rationale.}
Normalize is the boundary between surface authoring syntax and a regularized
AST.  It keeps later stages from duplicating parser-level sugar handling.

\paragraph{Formal view.}
Normalize is a total transformation over validated closed ASTs:
\[
  \mathsf{normalize}: A_c \rightarrow A_n .
\]
It preserves denotation while reducing the number of syntactic forms.  The
validator for \(A_n\) checks that eliminated sugar does not reappear.

\paragraph{Example: path and pipe sugar.}
The surface expression
\begin{verbatim}
x <- input |> NN.layernorm@ln_f eps=1e-5
\end{verbatim}
normalizes to an ordinary call whose path is an explicit argument:
\begin{verbatim}
x <- NN.layernorm @ln_f input eps=1e-5
\end{verbatim}
The exact rendered form may differ, but the AST no longer represents either a
pipe expression or callee-attached path sugar.

\paragraph{Example: loop protocol.}
A source loop with implied carry/yield is normalized so that the loop body has
an explicit final yield.  Flattening therefore sees one loop protocol rather
than several authoring conveniences.
